% ===== PRÉAMBULE CANONIQUE — à coller VERBATIM en tête de CHAQUE .tex =====
\documentclass[11pt,a4paper]{article}
\usepackage[utf8]{inputenc}\usepackage[T1]{fontenc}\usepackage{lmodern}
\usepackage[french]{babel}
\usepackage{amsmath,amssymb,mathtools}\usepackage{icomma}
\usepackage[version=4]{mhchem}          % \ce{...} : équations & formules
\usepackage{siunitx}\sisetup{locale=FR} % \qty{}{}, \si{}, \ang{}
\usepackage{chemfig}                    % \chemfig{...} : structures développées
\usepackage{xcolor}\usepackage{colortbl}
\usepackage{tikz}\usepackage{pgfplots}\pgfplotsset{compat=1.18}
\usepackage[most]{tcolorbox}\usepackage{enumitem}\usepackage{array,booktabs}
\usepackage[a4paper,margin=2.2cm]{geometry}
\usetikzlibrary{arrows.meta,calc,angles,quotes,positioning,patterns,backgrounds,shapes.geometric}
\definecolor{primary}{RGB}{222,49,99}\definecolor{primarydark}{RGB}{150,22,55}
\definecolor{defcol}{RGB}{0,114,178}\definecolor{methcol}{RGB}{52,92,148}
\definecolor{excol}{RGB}{0,135,90}\definecolor{attncol}{RGB}{200,40,55}
\tcbset{boxbase/.style={enhanced,breakable,boxrule=0.9pt,arc=2.5pt,
  left=3mm,right=3mm,top=2.3mm,bottom=2.3mm,fonttitle=\bfseries,coltitle=white}}
% --- boîtes TUTORIEL (noms EXACTS, ne pas renommer) ---
\newtcolorbox{cle}[1][]{boxbase,colback=defcol!6,colframe=defcol,
  title={\textsc{À retenir}\ifx\relax#1\relax\relax\else\textmd{ : \itshape #1}\fi}}
\newtcolorbox{methode}[1][]{boxbase,colback=methcol!7,colframe=methcol,sharp corners,
  title={\textsc{Méthode}\ifx\relax#1\relax\relax\else\textmd{ --- \itshape #1}\fi}}
\newtcolorbox{exemplebox}[1][]{boxbase,colback=excol!5,colframe=excol,
  title={\textsc{Exemple}\ifx\relax#1\relax\relax\else\textmd{ : \itshape #1}\fi}}
\newtcolorbox{piege}[1][]{boxbase,colback=attncol!6,colframe=attncol,
  title={\textsc{Piège}\ifx\relax#1\relax\relax\else\textmd{ : \itshape #1}\fi}}
\newtcolorbox{remarque}[1][]{boxbase,colback=black!4,colframe=black!45,
  title={\textsc{Remarque}\ifx\relax#1\relax\relax\else\textmd{ : \itshape #1}\fi}}
% --- boîtes EXERCICE / CORRIGÉ (noms EXACTS, auto-comptées) ---
\newtcolorbox[auto counter]{exo}[1][]{boxbase,colback=primary!4,colframe=primary,
  title={\textsc{Exercice}~\thetcbcounter\ifx\relax#1\relax\relax\else\textmd{ --- \itshape #1}\fi}}
\newtcolorbox[auto counter]{corrige}[1][]{boxbase,colback=excol!4,colframe=excol,
  title={\textsc{Corrigé}~\thetcbcounter\ifx\relax#1\relax\relax\else\textmd{ --- \itshape #1}\fi}}
\newcommand{\R}{\mathbb{R}}\newcommand{\N}{\mathbb{N}}\newcommand{\Z}{\mathbb{Z}}\newcommand{\ds}{\displaystyle}
% (un \usepackage{fancyhdr}+en-tête est possible ; il est ignoré côté web)
% =========================================================================
\begin{document}
\begin{exo}[l'ion perchlorate \ce{ClO4-}]
L'ion perchlorate est un polluant toxique pour la thyroïde. Applique la méthode des $5$ étapes à
\ce{ClO4-} : $n_{\max}$, $n_v$, nombre de liaisons, nombre de doublets libres. Précise la
répartition des doublets libres.
\end{exo}

\begin{exo}[le nitrite d'hydrogène \ce{HNO2}]
La connectivité réelle est \ce{H-O-N=O} (le \ce{H} est porté par un oxygène, pas par l'azote).
\begin{enumerate}[label=\alph*)]
  \item Calcule $n_{\max}$ et $n_v$.
  \item Calcule le nombre de liaisons et de doublets libres.
  \item Décris la structure : combien de doublets libres sur l'azote, sur chaque oxygène ?
\end{enumerate}
\end{exo}

\begin{exo}[l'ion sulfate \ce{SO4^2-}]
Applique la méthode complète à \ce{SO4^2-} (atome central \ce{S}). Donne $n_{\max}$, $n_v$, le
nombre de liaisons et de doublets libres, puis compare tes résultats à ceux de \ce{ClO4-} de
l'exercice~1 : que remarques-tu ?
\end{exo}

\begin{exo}[une molécule radicalaire : \ce{NO}]
Le monoxyde d'azote \ce{NO} possède un nombre \emph{impair} d'électrons.
\begin{enumerate}[label=\alph*)]
  \item Calcule $n_{\max}$ et $n_v$. Pourquoi doit-on utiliser la correction radicalaire ?
  \item Applique $n_l = n_{\max} - n_v - 1$, puis donne le nombre de liaisons.
  \item Calcule $n_{nl}$ et interprète le résultat (doublets libres et électron célibataire).
\end{enumerate}
\end{exo}

\begin{exo}[quand la méthode atteint sa limite : \ce{SF6}]
Applique la méthode des $5$ étapes à \ce{SF6}.
\begin{enumerate}[label=\alph*)]
  \item Calcule $n_{\max}$, $n_v$ et le nombre de liaisons prédit par la méthode.
  \item L'expérience montre que le soufre forme en réalité $6$ liaisons \ce{S-F}. Explique
        pourquoi la méthode simplifiée échoue ici et pourquoi la réalité est permise.
\end{enumerate}
\end{exo}

\end{document}
